\documentclass[10pt,a4paper]{article}
\usepackage{amsmath,amssymb,amsthm,graphicx}
% Card-local minimal preamble (journal unit sigma-2023-023-spectral-asymmetry-and-index-theory-on-manif, P2-ahl16 precedent).
%% Verbatim from the pinned source: the author theorem environments and macro block,
%% so every macro and environment of the counted spans is available.
% Omitted and why: the documentclass line (the journal class sigma is not on the host;
%% the assembler supplies the class) and page-geometry/font blocks (they fight the
%% forced xelatex engine).
\usepackage{amsmath,amssymb,amsthm,mathtools}


\numberwithin{equation}{section}
\newtheorem{Theorem}{Theorem}[section]
\newtheorem{Corollary}[Theorem]{Corollary}
\newtheorem{Lemma}[Theorem]{Lemma}
\newtheorem{Proposition}[Theorem]{Proposition}
 { \theoremstyle{definition}
\newtheorem{Definition}[Theorem]{Definition}
\newtheorem{Note}[Theorem]{Note}
\newtheorem{Example}[Theorem]{Example}
\newtheorem{Remark}[Theorem]{Remark} }
 \DeclareMathOperator{\Tr}{Tr}
 \DeclareMathOperator{\tr}{tr}
 \DeclareMathOperator{\Str}{Str}
 \DeclareMathOperator{\id}{id}
\DeclareMathOperator{\im}{im}
\DeclareMathOperator{\spec}{spec}
\DeclareMathOperator{\AS}{AS}
\DeclareMathOperator{\HS}{HS}
\DeclareMathOperator{\APS}{APS}
\DeclareMathOperator{\Dir}{Dir}
\DeclareMathOperator{\Neu}{Neu}
\DeclareMathOperator{\Id}{Id}
\DeclareMathOperator{\dom}{dom}
\DeclareMathOperator{\Ei}{Ei}
\DeclareMathOperator{\Idem}{Idem}
 \DeclareMathOperator{\Stab}{Stab}
\DeclareMathOperator{\Ad}{Ad}
\DeclareMathOperator{\ad}{ad}
\DeclareMathOperator{\ind}{index}
\DeclareMathOperator{\ext}{ext}
\DeclareMathOperator{\grad}{grad}
\DeclareMathOperator{\End}{End}
\DeclareMathOperator{\Hom}{Hom}
\DeclareMathOperator{\sgn}{sgn}
\DeclareMathOperator{\ch}{ch}
\DeclareMathOperator{\Todd}{Todd}
\DeclareMathOperator{\topp}{top}
\DeclareMathOperator{\reg}{reg}
\DeclareMathOperator{\mat}{mat}
\DeclareMathOperator{\diag}{diag}
\DeclareMathOperator{\Lie}{Lie}
\DeclareMathOperator{\rank}{rank}
\DeclareMathOperator{\sign}{sign}
\DeclareMathOperator{\diam}{diam}
\DeclareMathOperator{\red}{red}
\DeclareMathOperator{\LIM}{LIM}
\DeclareMathOperator{\Imag}{Im}
\DeclareMathOperator{\Real}{Re}
\DeclareMathOperator{\Spin}{Spin}
\DeclareMathOperator{\OO}{O}
\DeclareMathOperator{\SO}{SO}
\DeclareMathOperator{\SL}{SL}
\DeclareMathOperator{\GL}{GL}
\DeclareMathOperator{\SU}{SU}
\DeclareMathOperator{\U}{U}
\DeclareMathOperator{\Cl}{Cl}
 \DeclareMathOperator{\Int}{int}
 \DeclareMathOperator{\Ind}{Ind}
 \DeclareMathOperator{\ev}{ev}
 \DeclareMathOperator{\Res}{Res}
 \DeclareMathOperator{\Thom}{Thom}
 \DeclareMathOperator{\Span}{span}
 \DeclareMathOperator{\supp}{supp}
\DeclareMathOperator{\pt}{pt}
 \DeclareMathOperator{\even}{even}
 \DeclareMathOperator{\odd}{odd}
 \DeclareMathOperator{\DInd}{D-Ind}
\DeclareMathOperator{\KInd}{K-Ind}
\DeclareMathOperator{\vol}{vol}
\newcommand{\Spinc}{\Spin^c}
 \newcommand{\R}{\mathbb{R}}
 \newcommand{\C}{\mathbb{C}}
 \newcommand{\N}{\mathbb{N}}
 \newcommand{\Z}{\mathbb{Z}}
 \newcommand{\Q}{\mathbb{Q}}
 %\newcommand{\K}{\mathbb{K}}
 \newcommand{\bT}{\mathbb{T}}
 \newcommand{\bP}{\mathbb{P}}
\newcommand{\g}{\mathfrak{g}}
\newcommand{\h}{\mathfrak{h}}
\newcommand{\p}{\mathfrak{p}}
\newcommand{\kg}{\mathfrak{g}}
\newcommand{\kt}{\mathfrak{t}}
\newcommand{\ka}{\mathfrak{a}}
\newcommand{\XX}{\mathfrak{X}}
\newcommand{\kh}{\mathfrak{h}}
\newcommand{\kp}{\mathfrak{p}}
\newcommand{\kk}{\mathfrak{k}}
\newcommand{\kn}{\mathfrak{n}}
\newcommand{\cA}{\mathcal{A}}
\newcommand{\cE}{\mathcal{E}}
\newcommand{\calL}{\mathcal{L}}
\newcommand{\calH}{\mathcal{H}}
\newcommand{\cO}{\mathcal{O}}
\newcommand{\cB}{\mathcal{B}}
\newcommand{\cK}{\mathcal{K}}
\newcommand{\cP}{\mathcal{P}}
\newcommand{\cN}{\mathcal{N}}
\newcommand{\calD}{\mathcal{D}}
\newcommand{\cC}{\mathcal{C}}
\newcommand{\calS}{\mathcal{S}}
\newcommand{\cM}{\mathcal{M}}
\newcommand{\cU}{\mathcal{U}}
\newcommand{\cF}{\mathcal{F}}
\newcommand{\cCM}{\cC}
\newcommand{\PM}{P}
\newcommand{\DM}{D}
\newcommand{\LM}{L}
\newcommand{\vM}{v}
\newcommand{\sumGam}{\textstyle{\sum_{\Gamma}} }
\newcommand{\sigDg}{\sigma^D_g}
\newcommand{\Bigwedge}{\textstyle{\bigwedge}}
\newcommand{\ii}{\sqrt{-1}}
\newcommand{\Ubar}{\overline{U}}
\usepackage{cleveref}
\crefname{originalequation}{original equation}{original equations}
\Crefname{originalequation}{Original equation}{Original equations}
\makeatletter
\newlabel{eq D inv infty}{{3.1}{source}{}{sourcepointer-1}{}}
\newlabel{eq D inv infty@cref}{{[originalequation][0][]3.1}{[1][0][]source}{}{}{}}
\makeatother
\crefname{originalequation}{original equation}{original equations}
\Crefname{originalequation}{Original equation}{Original equations}
\makeatletter
\newlabel{eq D0+ P}{{3.20}{source}{}{sourcepointer-2}{}}
\newlabel{eq D0+ P@cref}{{[originalequation][0][]3.20}{[1][0][]source}{}{}{}}
\makeatother
\crefname{originalequation}{original equation}{original equations}
\Crefname{originalequation}{Original equation}{Original equations}
\makeatletter
\newlabel{eq D0- 1-P}{{3.21}{source}{}{sourcepointer-3}{}}
\newlabel{eq D0- 1-P@cref}{{[originalequation][0][]3.21}{[1][0][]source}{}{}{}}
\makeatother
\crefname{originalequation}{original equation}{original equations}
\Crefname{originalequation}{Original equation}{Original equations}
\makeatletter
\newlabel{eq DC+}{{3.9}{source}{}{sourcepointer-4}{}}
\newlabel{eq DC+@cref}{{[originalequation][0][]3.9}{[1][0][]source}{}{}{}}
\makeatother
\crefname{originalequation}{original equation}{original equations}
\Crefname{originalequation}{Original equation}{Original equations}
\makeatletter
\newlabel{eq DC-}{{3.10}{source}{}{sourcepointer-5}{}}
\newlabel{eq DC-@cref}{{[originalequation][0][]3.10}{[1][0][]source}{}{}{}}
\makeatother
\crefname{originalequation}{original equation}{original equations}
\Crefname{originalequation}{Original equation}{Original equations}
\makeatletter
\newlabel{eq S0++}{{3.14}{source}{}{sourcepointer-6}{}}
\newlabel{eq S0++@cref}{{[originalequation][0][]3.14}{[1][0][]source}{}{}{}}
\makeatother
\crefname{originalequation}{original equation}{original equations}
\Crefname{originalequation}{Original equation}{Original equations}
\makeatletter
\newlabel{eq def P}{{3.4}{source}{}{sourcepointer-7}{}}
\newlabel{eq def P@cref}{{[originalequation][0][]3.4}{[1][0][]source}{}{}{}}
\makeatother
\crefname{originallemma}{original lemma}{original lemmas}
\Crefname{originallemma}{Original lemma}{Original lemmas}
\makeatletter
\newlabel{lem D0pm}{{3.7}{source}{}{sourcepointer-8}{}}
\newlabel{lem D0pm@cref}{{[originallemma][0][]3.7}{[1][0][]source}{}{}{}}
\makeatother
\crefname{originallemma}{original lemma}{original lemmas}
\Crefname{originallemma}{Original lemma}{Original lemmas}
\makeatletter
\newlabel{lem T Dpsi}{{3.6}{source}{}{sourcepointer-9}{}}
\newlabel{lem T Dpsi@cref}{{[originallemma][0][]3.6}{[1][0][]source}{}{}{}}
\makeatother
\crefname{originallemma}{original lemma}{original lemmas}
\Crefname{originallemma}{Original lemma}{Original lemmas}
\makeatletter
\newlabel{lem T isom}{{3.5}{source}{}{sourcepointer-10}{}}
\newlabel{lem T isom@cref}{{[originallemma][0][]3.5}{[1][0][]source}{}{}{}}
\makeatother
\title{Invertibility at infinity and the correspondence of invertibility and self-adjointness between the two model operators}
\author{Hochs, Peter, Wang, Hang}
\date{doi:10.3842/SIGMA.2023.023; CC BY 4.0}
\begin{document}
\maketitle
\setcounter{section}{3}
\setcounter{equation}{22}
\setcounter{Theorem}{8}
\noindent\textit{Editorial note.} All delivered mathematical text is the author's own, copied from the pinned journal source file sources/sigma-2023-023-spectral-asymmetry-and-index-theory-on-manif/pinned-original-source.tex (Spectral Asymmetry and Index Theory on Manifolds with Generalised Hyperbolic Cusps); no statement or proof step was rewritten, reordered or removed.\par\medskip
\noindent\textit{Editorial note.} The counted proof is the author's complete argument: the counted statement is the author's (a)-(e) enumerate verbatim; the external references eq D inv infty, eq DC+, eq D0+ P, eq DC-, eq D0- 1-P, eq S0++, eq def P, lem T isom, lem T Dpsi and lem D0pm are navigation stubs carrying the article's printed numbers; the display label eq T W2 lies inside the span and is internal to the card.\par\medskip
\noindent\textit{Editorial note.} The label prop DC D0 names the counted statement, and the proof environment immediately following it in the source is delivered verbatim as its proof.\par\medskip
\noindent\textit{Editorial note.} The counted statement is carried by the author's own Proposition environment in the article source.\par\medskip
\noindent\textit{Editorial note.} Journal-unit assembly: card-local preamble (author theorem environments and macros verbatim), article class, no bibliography in the counted span (the span carries no \textbackslash{}cite commands); references leaving the span are declared as navigation stubs carrying the article's printed numbers.\par\medskip
\noindent\textit{Editorial note.} COUNTER REPAIR 2026-10-14: the r001 seeds were derived from the census v3.5+ printed field, which counts \textbackslash{}section and \textbackslash{}subsection jointly (subsection-inclusive) and diverges from the article's own numbering; the seeds above are recomputed by the lane's true-numbering rule (\{"section": 3, "equation": 22, "Theorem": 8\}).\par\medskip
\noindent\textit{Editorial scope.} the counted statement and its own proof environment only\par\medskip
\noindent\textit{Allowed standard prerequisites (editorial).}\begin{itemize}
\item The two first-order operators on the cylinder and the compact side. The operators D\_C\textasciicircum{}+/- on L\textasciicircum{}2(S\_C\textasciicircum{}+/-) with domains W\textasciicircum{}2\_D(S\_C\textasciicircum{}+/-; P or 1-P) and the model operators tilde D\_C\textasciicircum{}+/- on L\textasciicircum{}2(tilde S\_C\textasciicircum{}+) with domains W\textasciicircum{}2\_\{tilde D\_C\}(tilde S\_C\textasciicircum{}+; P or 1-P), together with the projection P of eq def P.. The pairs of operators the proposition relates; declared as the article's standing notation, no external stub beyond the operator-display stubs.
\end{itemize}\medskip
\ifdefined\hypertarget\hypertarget{sourcepointer-1}{}\fi
\noindent\textit{Source reference.} Original equation 3.1: sources/sigma-2023-023-spectral-asymmetry-and-index-theory-on-manif/pinned-original-source.tex, source lines 692--692. Navigation pointer invoked at line(s) 1014: the counted text cites that numbered equation of the same article, which is not retained in this card, so the reference is declared as a pointer carrying the number the article itself prints for it (the section and the place of the equation inside that section, derived by the census v3.1 rule for the section-scoped equation counter); the cited equation is an allowed input and none of it is reproduced inside the span.\par\medskip
\ifdefined\hypertarget\hypertarget{sourcepointer-2}{}\fi
\noindent\textit{Source reference.} Original equation 3.20: sources/sigma-2023-023-spectral-asymmetry-and-index-theory-on-manif/pinned-original-source.tex, source lines 1005--1005. Navigation pointer invoked at line(s) 1018: the counted text cites that numbered equation of the same article, which is not retained in this card, so the reference is declared as a pointer carrying the number the article itself prints for it (the section and the place of the equation inside that section, derived by the census v3.1 rule for the section-scoped equation counter); the cited equation is an allowed input and none of it is reproduced inside the span.\par\medskip
\ifdefined\hypertarget\hypertarget{sourcepointer-3}{}\fi
\noindent\textit{Source reference.} Original equation 3.21: sources/sigma-2023-023-spectral-asymmetry-and-index-theory-on-manif/pinned-original-source.tex, source lines 1006--1006. Navigation pointer invoked at line(s) 1019: the counted text cites that numbered equation of the same article, which is not retained in this card, so the reference is declared as a pointer carrying the number the article itself prints for it (the section and the place of the equation inside that section, derived by the census v3.1 rule for the section-scoped equation counter); the cited equation is an allowed input and none of it is reproduced inside the span.\par\medskip
\ifdefined\hypertarget\hypertarget{sourcepointer-4}{}\fi
\noindent\textit{Source reference.} Original equation 3.9: sources/sigma-2023-023-spectral-asymmetry-and-index-theory-on-manif/pinned-original-source.tex, source lines 772--772. Navigation pointer invoked at line(s) 1018: the counted text cites that numbered equation of the same article, which is not retained in this card, so the reference is declared as a pointer carrying the number the article itself prints for it (the section and the place of the equation inside that section, derived by the census v3.1 rule for the section-scoped equation counter); the cited equation is an allowed input and none of it is reproduced inside the span.\par\medskip
\ifdefined\hypertarget\hypertarget{sourcepointer-5}{}\fi
\noindent\textit{Source reference.} Original equation 3.10: sources/sigma-2023-023-spectral-asymmetry-and-index-theory-on-manif/pinned-original-source.tex, source lines 776--776. Navigation pointer invoked at line(s) 1019: the counted text cites that numbered equation of the same article, which is not retained in this card, so the reference is declared as a pointer carrying the number the article itself prints for it (the section and the place of the equation inside that section, derived by the census v3.1 rule for the section-scoped equation counter); the cited equation is an allowed input and none of it is reproduced inside the span.\par\medskip
\ifdefined\hypertarget\hypertarget{sourcepointer-6}{}\fi
\noindent\textit{Source reference.} Original equation 3.14: sources/sigma-2023-023-spectral-asymmetry-and-index-theory-on-manif/pinned-original-source.tex, source lines 890--890. Navigation pointer invoked at line(s) 1037: the counted text cites that numbered equation of the same article, which is not retained in this card, so the reference is declared as a pointer carrying the number the article itself prints for it (the section and the place of the equation inside that section, derived by the census v3.1 rule for the section-scoped equation counter); the cited equation is an allowed input and none of it is reproduced inside the span.\par\medskip
\ifdefined\hypertarget\hypertarget{sourcepointer-7}{}\fi
\noindent\textit{Source reference.} Original equation 3.4: sources/sigma-2023-023-spectral-asymmetry-and-index-theory-on-manif/pinned-original-source.tex, source lines 749--749. Navigation pointer invoked at line(s) 1046: the counted text cites that numbered equation of the same article, which is not retained in this card, so the reference is declared as a pointer carrying the number the article itself prints for it (the section and the place of the equation inside that section, derived by the census v3.1 rule for the section-scoped equation counter); the cited equation is an allowed input and none of it is reproduced inside the span.\par\medskip
\ifdefined\hypertarget\hypertarget{sourcepointer-8}{}\fi
\noindent\textit{Source reference.} Original lemma 3.7: sources/sigma-2023-023-spectral-asymmetry-and-index-theory-on-manif/pinned-original-source.tex, source lines 903--903. Navigation pointer invoked at line(s) 1052: the counted statement is phrased in notation from that statement of the same article, which is not retained in this card, so the reference is declared as a pointer carrying the number the article itself prints for it (the section and the place of the statement inside that section, derived by the census v3.1 rule for the shared theorem counter); the cited statement is an allowed input and none of it is reproduced inside the span.\par\medskip
\ifdefined\hypertarget\hypertarget{sourcepointer-9}{}\fi
\noindent\textit{Source reference.} Original lemma 3.6: sources/sigma-2023-023-spectral-asymmetry-and-index-theory-on-manif/pinned-original-source.tex, source lines 874--874. Navigation pointer invoked at line(s) 1031, 1052: the counted statement is phrased in notation from that statement of the same article, which is not retained in this card, so the reference is declared as a pointer carrying the number the article itself prints for it (the section and the place of the statement inside that section, derived by the census v3.1 rule for the shared theorem counter); the cited statement is an allowed input and none of it is reproduced inside the span.\par\medskip
\ifdefined\hypertarget\hypertarget{sourcepointer-10}{}\fi
\noindent\textit{Source reference.} Original lemma 3.5: sources/sigma-2023-023-spectral-asymmetry-and-index-theory-on-manif/pinned-original-source.tex, source lines 852--852. Navigation pointer invoked at line(s) 1031, 1052: the counted statement is phrased in notation from that statement of the same article, which is not retained in this card, so the reference is declared as a pointer carrying the number the article itself prints for it (the section and the place of the statement inside that section, derived by the census v3.1 rule for the shared theorem counter); the cited statement is an allowed input and none of it is reproduced inside the span.\par\medskip
\begin{Proposition}\label{prop DC D0}\quad
\begin{enumerate}\itemsep=0pt
\item[$(a)$] The operator $D$ is invertible at infinity in the sense of~\eqref{eq D inv infty} if and only if there are $b_0>0$ and $u\geq a$ such that for all $s \in \Gamma^{\infty}_c\big(\widetilde{S}_C\big)$ supported in $(u, \infty)$,
\[
\big\|\widetilde{D}_Cs\big\|_{L^2(\widetilde{S}_C, \vol_{B_N} \otimes {\rm d}y)} \geq b_0 \|s\|_{L^2(\widetilde{S}_C, \vol_{B_N} \otimes {\rm d}y)}.
\]
\item[$(b)$] The operator~\eqref{eq DC+} is invertible if and only if the operator~\eqref{eq D0+ P} is.
\item[$(c)$] The operator~\eqref{eq DC-} is invertible if and only if the operator~\eqref{eq D0- 1-P} is.
\item[$(d)$] The operator $D^-_C D_C^+$ on $L^2\big(S_C^+\big)$, with domain $W^2_D\big(S_C^+; P\big)$, is self-adjoint
 if and only if the operator $\widetilde{D}_C^- \widetilde{D}_C^+$ on $L^2\big(\widetilde{S}_C^+\big)$, with domain $W^2_{\widetilde{D}_C}\big(\widetilde{S}_C^+; P\big)$ is.
\item[$(e)$] The operator
the operator $D^+_C D_C^-$ on $L^2\big(S_C^-\big)$, with domain $W^2_D\big(S_C^-; 1-P\big)$
 is self-adjoint if and only if the operator
 $\widetilde{D}_C^+ \widetilde{D}_C^-$ on $L^2\big(\widetilde{S}_C^+\big)$, with domain $W^2_{\widetilde{D}_C}\big(\widetilde{S}_C^+; 1-P\big)$
 is.
\end{enumerate}
\end{Proposition}

\begin{proof}
Part (a) follows from Lemmas~\ref{lem T isom} and~\ref{lem T Dpsi}.

For any $s \in \Gamma^{\infty}\big(\widetilde{S}_C\big)$, we have
\[
(Ts)|_N = \Phi(a) s|_N,
\]
because we view $N$ as embedded into $\overline{C}$ as $N\times \{a\}$ and into $N \times [0,\infty)$ as $N \times \{0\}$. Therefore, for any such section, $P(s|_N) = 0$ if and only if $P (Ts|_N) = 0$, and similarly for $1-P$. Hence the operator $T$, together with~\eqref{eq S0++}, defines unitary isomorphisms
\begin{align}\label{eq T W2}
\begin{split}
&T\colon\ W^1_{\widetilde{D}_C}\big(\widetilde{S}_C^+; P\big) \xrightarrow{\cong} W^1_{D}\big(S_C^+; P\big), \\
&T \circ \tau \colon\ W^1_{\widetilde{D}_C}\big(\widetilde{S}_C^+;1- P\big) \xrightarrow{\cong}W^1_{D}\big(S_C^-;1- P\big), \\
&T\colon\ W^2_{\widetilde{D}_C}\big(\widetilde{S}_C^+; P\big) \xrightarrow{\cong}W^2_{D}\big(S_C^+; P\big),\\
&T\circ \tau \colon\ W^2_{\widetilde{D}_C}\big(\widetilde{S}_C^+; 1-P\big) \xrightarrow{\cong}W^2_{D}\big(S_C^-; 1-P\big).
\end{split}
\end{align}
In the second and fourth lines, we use the fact that by definition~\eqref{eq def P} of $P$, we have
\[
P|_{L^2(S^-|_N)} = \tau_+ P^+ \tau_+^{-1},
\]
so that the isomorphisms $T \circ \tau$ preserve the given boundary conditions.

Under the isomorphisms~\eqref{eq T W2}, the pairs of operators in parts (b)--(e) correspond to each other. Here we again used Lemmas~\ref{lem T isom} and~\ref{lem T Dpsi}, and also Lemma~\ref{lem D0pm}.
\end{proof}

\end{document}
